% ============================================================
% 04-experiments.tex
% ============================================================
\section{Experimental Results}
\label{sec:exp}

This section presents the empirical evaluation of \system{}, detailing the benchmark suites, infrastructure, evaluation metrics, three experimental scenarios, discussion answering each research question, and threats to validity.

%------------------------------------------------
\subsection{Experimental Benchmark and Protocols}
\label{sec:dataset}

To ensure rigorous validation and avoid data contamination, our empirical evaluation is established upon a corpus of $200$ curated administrative queries grounded in official Can Tho University (CTU) statutes. Following standard benchmark protocol, this corpus is partitioned into: (1)~\textbf{Development Benchmark ($N=150$ queries):} Spanning 9 institutional categories (\texttt{actual\_tuition}, \texttt{academic\_rules}, \texttt{scholarship}, \texttt{student\_loan}, \texttt{other}, \texttt{social\_support}, \texttt{academic\_program}, \texttt{exemption\_policy}, \texttt{exemption\_basis}), used for offline index calibration, prompt engineering, and qualitative failure analysis; and (2)~\textbf{Held-Out Test Set ($N=50$ queries):} A strictly separated, out-of-distribution evaluation suite covering 6 operational lanes (40\% tuition, 20\% academic regulations, 10\% scholarships, 10\% student credit loans, 10\% tuition exemptions, and 10\% administrative procedures). All reported primary benchmark metrics in Tables~\ref{tab:retrieval_evaluation} and~\ref{tab:e2e_evaluation} are evaluated strictly on this blind Held-Out Test Set.

\noindent\textbf{Execution and Evaluation Invariants:} Generative and automated evaluation tasks were executed using Google Gemini 2.5 Flash Lite (\texttt{gemini-2.5-flash-lite}) via Google Cloud Vertex AI under strict zero-temperature decoding ($\text{temperature}=0.0$) to guarantee determinism. To eliminate stochastic variance, each query was evaluated across $R=3$ independent repetitions, yielding $1,050$ end-to-end generation and Ragas evaluation turns ($7\text{ configurations} \times 50\text{ queries} \times 3\text{ reps}$).

\noindent\textbf{Computational Infrastructure:} Neural cross-encoder reranking was executed on an NVIDIA GeForce RTX 4090 GPU (24~GB VRAM). Relational curricula are hosted in Neo4j 5 Enterprise with APOC extensions; document parent chunks (2,800 characters, 100 overlap) reside in PostgreSQL, while child chunks (400 characters, 50 overlap) are indexed in Qdrant with \texttt{bkai-foundation-models/vietnamese-bi-encoder} (768-d). Neural candidate reranking uses \texttt{BAAI/bge-reranker-v2-m3}. Multi-agent coordination was orchestrated via LangGraph under a StateGraph supervisor with Redis 7 caching.

%------------------------------------------------
\subsection{Evaluation Metrics}
\label{sec:metrics}

\textbf{Multi-Agent Coordination Metrics (Scenario 3):} Following isolated evaluation protocols, we assess: (1)~\textbf{Agent Accuracy \& Macro-F1:} proportion of queries routed to the target specialist alongside macro-averaged F1 across the four domains; (2)~\textbf{Intent Accuracy:} accuracy of fine-grained intent classification; (3)~\textbf{Tool Selection, Argument Exact Match (EM), and Result Accuracy:} whether the correct execution path is invoked, whether extracted arguments match ground truth with exact numeric and string equality, and whether outputs are mathematically exact; (4)~\textbf{End-to-End Pass Rate:} joint success across tool selection, argument binding, and execution output; (5)~\textbf{Robustness \& Tool-Gating Accuracy:} correct tool triggering on complete inputs versus deliberate tool suppression on malformed or policy-only queries; (6)~\textbf{Safe Result Behavior:} explicit abstention or clarification without parameter fabrication; and (7)~\textbf{Bounded Completion Rate:} proportion of decisions terminating within operational bounds ($\text{max\_tool\_calls}=1, \text{timeout}=60.0\,\text{s}$).

\textbf{Retrieval Quality Metrics (Scenario 1):} Evaluated over top-$k$ returned passages against gold document identifiers: (1)~\textbf{Hit@$k$ ($k \in \{1, 3\}$):} proportion of queries where at least one gold document is retrieved within top-$k$; (2)~\textbf{Precision@5 (P@5)} and \textbf{Recall@5 (R@5):} proportion of relevant passages in top-5, and proportion of gold passages retrieved; and (3)~\textbf{Mean Reciprocal Rank (MRR@10):} reciprocal rank of the first relevant document.

\textbf{End-to-End Generation Metrics (Scenario 2):} Evaluated using Ragas~\cite{es2023ragas} and ground-truth alignment: (1)~\textbf{Answer Relevancy (AR):} whether the generated answer directly addresses the question; (2)~\textbf{Context Recall (CR):} proportion of factual information required by the reference answer present in the retrieved Top-7 context; (3)~\textbf{Context Precision (CP):} concentration of relevant passages at higher ranks in Top-7; (4)~\textbf{Answer Correctness (AC):} factual agreement between generated answer and ground truth; (5)~\textbf{Factual Exact Match (Fact EM):} strict numerical and cohort exact matching on unseen queries; and (6)~\textbf{Source AP:} average precision of cited regulatory sources.

%------------------------------------------------
\subsection{Experimental Scenarios and Benchmark Results}
\label{sec:scenarios}

Table~\ref{tab:rq_mapping} maps each research gap and question to its corresponding architectural component, experiment, and empirical evidence.

\begin{table}[!htbp]
\centering
\caption{Systematic mapping between Research Gaps, RQs, Proposed Modules, and Target Metrics.}
\label{tab:rq_mapping}
\resizebox{\columnwidth}{!}{
\begin{tabular}{lllll}
\toprule
\textbf{Gap \& RQ} & \textbf{Proposed Architectural Component} & \textbf{Experiment} & \textbf{Primary Metrics} & \textbf{Key Empirical Evidence} \\
\midrule
\textbf{RQ1 (Multi-Agent)} & Supervisor + 4 Domain Specialists (C1) & Scenario 3 & Routing Acc., Tool EM, Safety & 97.0\% Routing Acc., 100\% Bounded \\
\textbf{RQ2 (Allocation)}  & Neo4j Graph vs. Hybrid RAG (C2)        & Scenario 2 & Context Precision, Ans. Correctness & T4 top Fact EM (62.03\%), CP (0.6977) \\
\textbf{RQ3 (Routing)}     & Subspace Governance + Reranking (C3)   & Scenario 1 & Hit@$k$, Recall@5, MRR@10            & E5 attains 0.9600 Hit@1, 0.9700 MRR@10 \\
\bottomrule
\end{tabular}
}
\end{table}

\subsubsection{Scenario 1 --- Retrieval Component Stacking (E1--E5):}
\emph{Objective \& RQ Alignment:} Scenario~1 investigates \textbf{RQ3} and supports \textbf{RQ2} by evaluating the empirical contribution of progressively stacking retrieval components over the 50-query Held-Out Test Set.
\emph{Configurations:} We compare five cumulative setups: (1)~\textbf{E1 (BM25 Lexical Baseline):} BM25 search over tokenized child chunks; (2)~\textbf{E2 (Dense Vector Semantic):} 768-d \texttt{vietnamese-bi-encoder} vector search in Qdrant; (3)~\textbf{E3 (Hybrid RRF):} BM25 and Dense fused via Reciprocal Rank Fusion ($k=60$); (4)~\textbf{E4 (Hybrid + Neural Reranker):} E3 augmented with \texttt{bge-reranker-v2-m3} cross-attention; and (5)~\textbf{E5 (CTU-Chat Proposed Retrieval):} integrating IG-SRR subspace routing, Neo4j relational grounding, and parent-child chunk mapping.

\emph{Results \& Interpretation:} As shown in Table~\ref{tab:retrieval_evaluation}, BM25 achieves a strong initial lexical baseline ($0.9000$ Hit@1 and $0.9333$ MRR@10) due to exact cohort codes and decree identifiers. Dense retrieval alone exhibits marked degradation ($0.3800$ Hit@1, $0.5237$ MRR@10) because vector embeddings conflate regulations across cohorts. Unreranked hybrid fusion (E3) balances lexical and semantic representations ($0.6800$ Hit@1, $0.7912$ MRR@10). Introducing cross-encoder reranking (E4) substantially improves candidate reordering ($0.8800$ Hit@1, $0.9267$ MRR@10, and $0.9800$ Hit@3). Finally, full heterogeneous integration with IG-SRR subspace routing and Neo4j relational grounding (E5) achieves the highest performance across all primary metrics: \textbf{0.9600 Hit@1} (48 out of 50 queries retrieving gold passages at rank 1), \textbf{0.9800 Hit@3}, and \textbf{0.9700 MRR@10}, outperforming BM25 by $+6.0\%$ in top-1 accuracy and verifying robust generalization on unseen queries.

\begin{table}[!htbp]
\centering
\caption{Scenario 1: Progressive retrieval component stacking on the Held-Out Test Set ($N=50$).}
\label{tab:retrieval_evaluation}
\resizebox{\columnwidth}{!}{
\begin{tabular}{lccccc}
\toprule
\textbf{Retrieval Configuration} & \textbf{Hit@1} & \textbf{Hit@3} & \textbf{P@5} & \textbf{Recall@5} & \textbf{MRR@10} \\
\midrule
E1: BM25 Lexical Baseline        & 0.9000 & 0.9400 & \textbf{0.2320} & \textbf{0.9667} & 0.9333 \\
E2: Dense Vector Semantic        & 0.3800 & 0.6600 & 0.1560 & 0.6433 & 0.5237 \\
E3: Hybrid RRF (BM25 + Dense)    & 0.6800 & 0.9200 & 0.2160 & 0.9133 & 0.7912 \\
E4: Hybrid + Neural Reranker     & 0.8800 & \textbf{0.9800} & 0.2200 & 0.9433 & 0.9267 \\
\textbf{E5: CTU-Chat Proposed Retrieval} & \textbf{0.9600} & \textbf{0.9800} & 0.2200 & 0.9433 & \textbf{0.9700} \\
\bottomrule
\end{tabular}
}
\vspace{0.5em}

\caption{Scenario 2: End-to-end QA generation and modular ablation on Held-Out Test Set ($N=50, R=3$, 1,050 total evaluation turns, Top-7 context).}
\label{tab:e2e_evaluation}
\resizebox{\columnwidth}{!}{
\begin{tabular}{llcccccc}
\toprule
\textbf{Cfg.} & \textbf{Ablation Variant Description} & \textbf{Fact EM} & \textbf{CP} & \textbf{CR} & \textbf{AR} & \textbf{AC} & \textbf{Source AP} \\
\midrule
T1 & BM25 Lexical Baseline                       & 0.5572 & 0.6674 & 0.8800 & 0.6151 & \textbf{0.6735} & 0.8954 \\
T2 & Dense Semantic Baseline                     & 0.5007 & 0.2770 & 0.3400 & 0.2809 & 0.3841 & 0.4651 \\
T3 & Hybrid RRF Baseline                         & 0.5451 & 0.4339 & 0.7733 & 0.5754 & 0.6016 & 0.7490 \\
\textbf{T4} & \textbf{CTU-Chat Full Stack (Proposed)}    & \textbf{0.6203} & \textbf{0.6977} & \textbf{0.8933} & \textbf{0.6195} & 0.6637 & \textbf{0.9311} \\
T5 & w/o Cross Reranker (\texttt{bge-m3})        & 0.5823 & 0.6297 & 0.8133 & 0.5263 & 0.6136 & 0.8168 \\
T6 & w/o Neo4j Knowledge Graph Grounding         & 0.5711 & 0.5683 & 0.8600 & 0.6173 & 0.6582 & 0.9011 \\
T7 & w/o Subspace Governance (Ungoverned)        & 0.6089 & 0.6947 & 0.8600 & 0.5815 & 0.6430 & \textbf{0.9311} \\
\bottomrule
\end{tabular}
}
\smallskip
\parbox{\columnwidth}{\footnotesize\emph{Metrics:} Fact EM: Factual Exact Match on tuition/credits; CP: Context Precision; CR: Context Recall; AR: Answer Relevancy; AC: Answer Correctness; Source AP: Source Average Precision. Normalized to $[0,1]$; higher is better.}
\end{table}

%------------------------------------------------
\subsubsection{Scenario 2 --- End-to-End QA Generation and Modular Ablation (T1--T7):}
\emph{Objective \& RQ Alignment:} Scenario~2 addresses \textbf{RQ2} and \textbf{RQ3} by evaluating end-to-end answer synthesis and modular ablation across seven configurations on the Held-Out Test Set ($1,050$ total evaluation turns across 3 independent repetitions).
\emph{Configurations:} (T1) BM25+LLM; (T2) Dense+LLM; (T3) Hybrid+LLM; (T4) Proposed \system{} Full Stack; (T5) w/o Reranker; (T6) w/o Neo4j Knowledge Graph; and (T7) w/o Governance Filter (removing subspace routing).
\emph{Results \& Interpretation:} As reported in Table~\ref{tab:e2e_evaluation}, the proposed full system (T4) decisively outperforms all baselines in factual fidelity: T4 attains the top \textbf{Factual Exact Match (Fact EM: 62.03\% vs. 55.72\% for BM25, a +6.31\% absolute gain)}, demonstrating superior precision on complex numeric tuition schedules and credit rules. In addition, T4 leads in Context Precision ($0.6977$ vs. $0.6674$ for BM25), Context Recall ($0.8933$ vs. $0.8800$), Source AP ($0.9311$ vs. $0.8954$), and Answer Relevancy ($0.6195$ vs. $0.6151$). While BM25 (T1) scores slightly higher on lexical Answer Correctness ($0.6735$) due to direct keyword overlap, T4 achieves superior structural grounding and eliminates hallucinations. Ablation results confirm the vital contribution of each module: omitting the cross-reranker (T5) reduces AR from $0.6195$ to $0.5263$ ($-0.0931$) and CR to $0.8133$; removing subspace governance (T7) degrades Answer Relevancy to $0.5815$ and Correctness to $0.6430$; and removing Knowledge Graph grounding (T6) causes a dramatic plunge in Context Precision from $0.6977$ to $0.5683$ ($-12.75\%$), proving its indispensable structural value.

%------------------------------------------------
\subsubsection{Scenario 3 --- Multi-Agent Routing, Tool Reliability, and Robustness:}
\emph{Objective \& RQ Alignment:} Scenario~3 directly addresses \textbf{RQ1} by evaluating multi-agent coordination, tool execution reliability, and bounded safety independently from retrieval or generation.
\emph{Experimental Setup:} Evaluated across three suites ($N=550$ runs): (1)~\textbf{Supervisor Routing (Panel A):} 100 single-domain benchmark queries ($N=300$ for LLM Supervisor vs. $N=100$ for Rule-Based Router); (2)~\textbf{Specialist Tool Reliability (Panel B):} 60 production cases ($N=180$ across 3 reps) covering 11 production tools and no-tool gating; and (3)~\textbf{Adversarial Robustness (Panel C):} 20 failure stress tests ($N=60$ across 3 reps) challenging boundary invariants.

\begin{table}[!htbp]
\centering
\caption{Scenario 3: Multi-agent routing, specialist tool execution reliability, and robustness under a maximum of one tool call per decision ($\text{max\_tool\_calls}=1$).}
\label{tab:scenario3_results}
\resizebox{\columnwidth}{!}{
\begin{tabular}{lccccccc}
\toprule
\multicolumn{8}{l}{\textbf{Panel A: Supervisor Routing and Specialist Selection ($N=100$ queries; Rule: 1 rep / LLM: 3 reps, $N=300$)}} \\
\midrule
\textbf{Router Variant} & \textbf{Runs ($N$)} & \textbf{Agent Acc.} & \textbf{Macro-F1} & \textbf{Intent Acc.} & \textbf{Bounded} & \textbf{p50 (ms)} & \textbf{p95 (ms)} \\
\midrule
Rule-Based Router          & 100 & 85.0\% & 69.2\% & 69.0\% & 100.0\% & \textbf{0.0} & \textbf{0.1} \\
LLM Supervisor (\system{}) & 300 & \textbf{97.0\%} & \textbf{97.8\%} & \textbf{95.0\%} & \textbf{100.0\%} & 795.5 & 981.7 \\
\midrule
\multicolumn{8}{l}{\textbf{Panel B: Specialist Tool Execution Reliability ($N=15$ runs per function, 3 reps)}} \\
\midrule
\textbf{Target Function} & \textbf{Execution Type} & \textbf{Selection} & \textbf{Arg. EM} & \textbf{Result Acc.} & \textbf{E2E Pass} & \textbf{Bounded} & \textbf{Graph hit} \\
\midrule
\texttt{tra\_cuu\_nganh} & LLM Tool Call & 100.0\% & 100.0\% & 100.0\% & 100.0\% & 100.0\% & -- \\
\texttt{so\_sanh\_nganh} & LLM Tool Call & 100.0\% & 100.0\% & 100.0\% & 100.0\% & 100.0\% & -- \\
\texttt{tim\_nganh} & LLM Tool Call & 100.0\% & 100.0\% & 100.0\% & 100.0\% & 100.0\% & -- \\
\texttt{xem\_chuoi\_tien\_quyet} & LLM Tool Call & 100.0\% & 100.0\% & 100.0\% & 100.0\% & 100.0\% & -- \\
\texttt{mon\_chung\_giua\_nganh} & LLM Tool Call & 93.3\% & 93.3\% & 100.0\% & 93.3\% & 100.0\% & -- \\
\texttt{tim\_nganh\_co\_mon} & LLM Tool Call & 100.0\% & 100.0\% & 100.0\% & 100.0\% & 100.0\% & -- \\
\texttt{tra\_cuu\_hoc\_phi\_graph} & LLM Tool Call & 80.0\% & 80.0\% & 100.0\% & 80.0\% & 100.0\% & 80.0\% \\
\texttt{tra\_cuu\_co\_so\_mien\_giam\_graph} & LLM Tool Call & 80.0\% & 80.0\% & 73.3\% & 73.3\% & 93.3\% & 78.6\% \\
\texttt{tra\_cuu\_quy\_dinh\_hoc\_phi} & LLM Tool Call & 100.0\% & 100.0\% & 100.0\% & 100.0\% & 100.0\% & 40.0\% \\
\texttt{tinh\_toan\_hoc\_phi} & LLM Tool Call & 100.0\% & 80.0\% & 80.0\% & 80.0\% & 100.0\% & -- \\
\texttt{tinh\_tien\_hoc\_bong} & LLM Tool Call & 100.0\% & 100.0\% & 100.0\% & 100.0\% & 100.0\% & -- \\
\texttt{no\_tool} & No Tool Call & 100.0\% & 100.0\% & 100.0\% & 100.0\% & 100.0\% & -- \\
\midrule
\multicolumn{8}{l}{\textbf{Panel C: Robustness and Bounded Failure Handling ($N=60$ runs under malformed/adversarial inputs)}} \\
\midrule
\textbf{Adversarial Gate} & \textbf{Runs ($N$)} & \textbf{Agent Acc.} & \textbf{Tool Gate} & \textbf{Safe Result} & \textbf{E2E Pass} & \textbf{Bounded} & \textbf{Notes} \\
\midrule
Failure Gate Robustness    & 60 & 100.0\% & 100.0\% & 98.3\% & 98.3\% & 100.0\% & -- \\
\bottomrule
\end{tabular}
}
\smallskip
\parbox{\columnwidth}{\footnotesize\emph{Note:} \texttt{tuition\_lookup} evaluates deterministic catalog lookup. Scholarship and fee reduction evaluate typed LLM tool selection.}
\end{table}

\emph{Results \& Interpretation:}
(1)~\emph{Supervisor vs. Heuristic Routing:} The LLM Supervisor attains \textbf{97.0\% Agent Accuracy}, \textbf{97.8\% Macro-F1}, and \textbf{95.0\% Intent Accuracy}, outperforming the rule router ($85.0\%$, $69.2\%$, and $69.0\%$) with sub-second p95 latency ($981.7\,\text{ms}$, p50 $795.5\,\text{ms}$).
(2)~\emph{Specialist Tool Execution:} Five academic tools, tuition policy lookup, scholarship calculator, and no-tool gating achieved \textbf{100.0\%} end-to-end pass.
(3)~\emph{Robustness \& Boundedness:} Under adversarial stress testing (Panel~C), \system{} achieved \textbf{100.0\% Tool/No-tool decision accuracy}, \textbf{98.3\% Safe Result behavior}, and \textbf{100.0\% bounded completion} within operational limits ($\text{max\_tool\_calls}=1, \text{timeout}=60\,\text{s}$).

%------------------------------------------------
\subsection{Discussion}
\label{sec:discussion}

\textbf{Answering RQ1 (Domain-Specialized Multi-Agent Architecture):} Centralizing routing under a structured supervisor with asymmetric privileges and bounds ($\text{max\_tool\_calls}=1$, timeout $=60\,\text{s}$) produced bounded isolated decisions in the Scenario~3 runs: the LLM Supervisor achieves $97.0\%$ agent accuracy and $97.8\%$ Macro-F1 across 300 evaluations, outperforming rule-based routing ($85.0\%$ and $69.2\%$). Production-tool evaluations reached $100.0\%$ end-to-end pass on 8 of 12 functions, while adversarial testing demonstrated $100.0\%$ tool-gating accuracy, $98.3\%$ safe-result behavior, and $100.0\%$ bounded completion.

\textbf{Answering RQ2 (Heterogeneous Knowledge Allocation \& The Graph Ablation):} Decomposing counseling knowledge across structured relational curricula (Neo4j) and unstructured regulatory decrees (hybrid RAG) is strongly substantiated by our modular ablations. 
A key empirical question concerns the \textit{without-graph ablation} (T6): \textit{Why did T6 achieve competitive context scores on certain single-entity queries during development?}
Our qualitative error analysis reveals that for \textbf{single-entity factual lookups} (e.g., standard tuition for a general major), unstructured text paragraphs alone provide complete answers; injecting structured graph nodes introduces slight context overhead and token dispersion, marginally diluting automated precision scores. 
However, for \textbf{multi-hop relational queries} (e.g., prerequisite academic chains, cohort-specific tuition progression, and credit-reduction criteria), the Knowledge Graph is indispensable. 
This is conclusively proven by the blind Held-Out Test Set: when evaluated against complex unseen queries, omitting Neo4j (T6) causes Context Precision to \textbf{plummet from 0.6977 to 0.5683 (a severe $-12.75\%$ absolute drop, $-18.3\%$ relative)}, while reducing Factual Exact Match from $62.03\%$ to $57.11\%$ ($-4.92\%$) and Context Recall from $0.8933$ to $0.8600$. Thus, Knowledge Graph grounding provides essential structural anchors that prevent factual hallucination in complex reasoning paths.

\textbf{Answering RQ3 (Representation-Aware Evidence Routing):} Across cumulative retrieval configurations, top-1 retrieval rises from $0.9000$ (E1) to a record \textbf{0.9600 Hit@1} ($48/50$ queries correct at rank 1) and \textbf{0.9700 MRR@10} for E5 on the Held-Out Test Set. In end-to-end ablation, T4's governed Graph injection and reranking achieve the highest Factual Exact Match ($62.03\%$), Context Precision ($0.6977$), and Source AP ($0.9311$), whereas removing governance routing (T7) degrades Answer Relevancy to $0.5815$ and Correctness to $0.6430$. While neural-reranked configurations exhibit higher measured latency ($10.99\,\text{ms}$ for E1 vs. $18{,}144.98\,\text{ms}$ for E5), unreranked hybrid fusion (E3, $43.62\,\text{ms}$) represents a practical low-latency alternative when computational budgets are constrained.

%------------------------------------------------
\subsection{Threats to Validity and Limitations}
\label{sec:limitations}

While derived from Can Tho University, the administrative taxonomy generalizes across Vietnamese higher education; multi-institutional validation remains planned. In addition, cross-attention reranking latency motivates future model distillation. Curricula graphs currently rely on offline schemas, and automated judges penalize out-of-scope abstentions, underestimating practical correctness when refusing ungrounded generation. Finally, while Scenario~3 validates bounded execution and tool reliability at the isolated supervisor and specialist decision level, end-to-end fault injection across external infrastructure services (e.g., Neo4j connection drops or Qdrant timeouts) represents an orthogonal system-resilience evaluation reserved for future work.
